\SourceLeaf{633}
\begin{center}[20]\par\medskip[\textit{C. F. Gauss to W. Weber.}]\end{center}
My esteemed friend,

Ever since the year began, demands of many different kinds have kept taking up and fragmenting my time, and my health has been so ill suited to sustained work that I have had no opportunity to go through the short paper you kindly sent me two months ago. Only now have I managed a cursory reading. It has shown me that the subject belongs to the same inquiries I worked on at considerable length about 10 years ago (especially in 1834–1836, I mean). Reading \textit{your paper} alone would be insufficient for a thorough and complete judgment on it. I would first have to immerse myself again in my own work from that period. That would take all the longer because a trial search through my papers has so far turned up only a few bits and pieces, though much more probably survives, even if it is not fully in order.

If I may venture a judgment from memory, after being away from these subjects for several years, I would think that \textsc{Ampère}, were he still alive, would object firmly from the outset if you were to use the formula
\[
-\frac{\alpha\alpha'}{rr}ii'\sin\theta\sin\theta'\cos\varepsilon\tag*{(I)}
\]
to express his fundamental law. His law is quite different, as the formula
\[
-\frac{\alpha\alpha'}{rr}ii'\left(\tfrac12\cos\theta\cos\theta'+\sin\theta\sin\theta'\cos\varepsilon\right)\tag*{(II)}
\]
shows. I also doubt that \textsc{Ampère} would be satisfied by the additional note you mention in a later letter. You present the difference there by saying that \textsc{Ampère}’s formula is a \textit{more general} one, of the form \mbox{$-\dfrac{\alpha\alpha'}{rr}(F\cos\theta\cos\theta'+G\sin\theta\sin\theta'\cos\varepsilon)$}. In your account, \textsc{Ampère} derived $F=\tfrac12G$ from experiments, while you think you can claim the value $F=0$ with equal justification because his experiments lack precision. In any other case I would concede that a third person might put the disagreement between you and \textsc{Ampère} somewhat as follows.

\SourceLeaf{634}
Whether one regards this (as you do) as only a modification of \textsc{Ampère}’s law, or

whether it amounts (as I believe \textsc{Ampère} would have to see it) to completely overturning his fundamental formula and putting a substantially different one in its place

would, when it comes down to it, be little more than an idle dispute over words. As I said, I would gladly concede this in any other case, since no one can be \textit{in verbis facilior} than I am. Here, though, it is a question of survival. \textsc{Ampère}’s whole theory of the interchangeability of magnetism and galvanic currents stands entirely on the correctness of formula II and would be lost altogether if another were chosen in its place.

I cannot take issue with you when you say that \textsc{Ampère}’s experiments are less than conclusive, especially since I do not have his classic paper to hand and cannot remember at all how he conducted the experiments. Even if he were to acknowledge their shortcomings himself, though, I doubt he would grant anyone the right to adopt a wholly different formula (I), which would bring his entire theory down, until that other formula had been borne out by \textit{truly decisive} experiments. You must have misunderstood the reservations that, according to your second letter, I myself expressed. I became convinced early on, and have remained so, that the interchangeability mentioned above \textit{necessarily} calls for \textsc{Ampère}’s formula II and allows only formulas that become identical to it for a closed current, \textit{if the action is to take place along the straight line joining the two current elements}. If that condition is dropped, one can indeed choose any of countless other forms, though for a closed current they must always lead to the same final result as his formula. One might add that, since for those purposes we are only ever concerned with effects at measurable distances, nothing would prevent us from supposing that there might also be other terms in the formula that take effect only at distances too small to measure (as molecular attraction comes in alongside gravitation). This might remove the difficulty posed by the repulsion of two successive elements of the same current.

\SourceLeaf{635}
To avoid misunderstanding, let me also point out that formula II above can be written as
\[
-\frac{\alpha\alpha'}{rr}ii'\left(-\tfrac12\cos\theta\cos\theta'+\sin\theta\sin\theta'\cos\varepsilon\right)
\]
and that I do not know whether \textsc{Ampère} (whose memoir, as I said, I do not have to hand) used the first notation or the second. The two mean the same thing. One writes the first form when measuring the angles $\theta$, $\theta'$ from the same (bounded) straight line, taking that line in the opposite direction as the side of the second angle. One writes the other form when considering a straight line of unspecified length and using it in the same direction both times to measure the angles $\theta$, $\theta'$. Likewise, one can replace the $-$ sign in front of the whole formula with a $+$ sign if attraction is counted as a positive effect instead of repulsion.

Perhaps I can get a little further back into these matters, which now feel so unfamiliar to me, before you give me the pleasure of a visit at the end of April or the beginning of May, as you have led me to hope. I would certainly have made my investigations known long ago had the thing I regarded as their very keystone not been missing when I broke them off
\begin{center}\textit{Nil actum reputans si quid superesset agendum}\end{center}
namely, \textit{the derivation} of the additional forces (which are added to the mutual action of portions of electricity at rest when those portions are in relative motion) \textit{from an action that is not instantaneous}, but propagates over time (in a similar way to light). I had been unable to do this at the time. As far as I can recall, I left the investigation without entirely losing hope that it might succeed later, although I held the subjective conviction, if I remember rightly, that one would first have to form a constructible conception of \textit{how} the propagation takes place.
\begin{center}—\quad —\quad —\quad —\quad —\quad —\quad —\quad —\quad —\quad —\quad —\quad —\quad —\quad —\quad —\quad —\end{center}
\noindent\begin{minipage}{\linewidth}
With warm greetings to your siblings and to Professor \textsc{Möbius}

Göttingen, 19 March 1845.\hfill always yours
\begin{flushright}C. F. \textsc{Gauss}.\end{flushright}
\end{minipage}

\SourceLeaf{636}
\begin{center}[21.]\end{center}
\noindent Linear quantity $=$ resistance of a given wire $=r$\par
\noindent Time quantity $=t$\par
\noindent Velocity $=\dfrac rt$\par
\noindent Density $=\dfrac1{tt}=\dfrac p{r^3}$, the density of water is approximately $\dfrac1{15\,000\,000(1'')^2}$\par
\noindent Expansibility of the liquid $=\dfrac{rr}{t^4}=\dfrac p{rtt}$\par
\noindent Specific elasticity at a given temperature $=\dfrac{rr}{tt}$\par
\noindent Accelerating force $=\dfrac r{tt}$\par
\noindent Mass $=\dfrac{r^3}{tt}=p$\par
\noindent Pressure $=\dfrac{r^4}{t^4}=\dfrac{pr}{tt}$\par
\noindent Effect $=$ vis viva $=$ torque $=\dfrac{r^5}{t^4}=\dfrac{prr}{tt}$\par
\noindent Power $=\dfrac{r^5}{t^5}=\dfrac{prr}{t^3}$\par
\noindent Terrestrial magnetism $=\BareRoot{\dfrac p{rtt}}$\par
\noindent Free magnetism $=$ intensity of an entire current $=\BareRoot{\dfrac{pr^3}{tt}}$\par
\noindent Specific intensity of a galvanic current $=\BareRoot{\dfrac{pr}{tt}}$\par
\noindent Exciting force of copper: zinc $=\BareRoot{\dfrac{pr^3}{t^4}}$\par
\noindent Conducting power of a given metal $=\dfrac t{rr}$\par

